% ===== 问题1的模型建立与求解 =====
\subsection{问题1的模型建立与求解}

% ==== 建模 ====
% 结构：开篇引文 → 编号步骤逐步建模 → 总结桥接
% 公式密度：每个编号步骤至少 6-7 个公式
% 文字串联：公式间用"将式(X)代入式(Y)可得"过渡
% 变量说明：每个公式后紧跟"其中，$X$表示……"（自然段，非列表）

% ==== 数据预处理（如适用，5段式结构） ====
% 第1段：预处理原因（数据来源、格式、存在问题）
% 第2段：数据理解与提取（特征数量、样本规模）
% 第3段：预处理步骤（清洗→标准化→特征提取，含公式）
% 第4段：预处理结果（处理前后统计量对比表）
% 第5段：总结（数据已规整，为建模提供输入）

% ==== 算法选择（5段式结构） ====
% 第1段：问题本质定性——"问题1本质上是一个【类型】问题"
% 第2段：候选算法 + 评估维度
% 对比表：booktabs 三线表（算法 + 4个评估维度，含中英文缩写 + 优良中差评级）
% 第3段：分析对比结果——"XX算法在XX方面最优"
% 第4段：总结 + 引出——"综上所述，采用XX算法"

% ==== 求解 ====
% 算法参数设置 + 求解过程描述

% ==== 结果分析 ====
% 每图 ≥ 3 句分析（趋势 + 归因 + 决策）
% 使用 \cref{fig:xxx} 引用图片（自动生成"图1"）
% 使用 \cref{tab:xxx} 引用表格（自动生成"表1"）
% 图片来源：../求解/问题1/图片/xxx.png
